\section*{الفصل \ref{ch:g7:volcanoes} --- البراكين}

\begin{solution}{exo:g7:volcanoes:1}
(أ) \omterm{def:g7:volcanoes:vent}{الفوهة}؛ (ب) \omterm{def:g7:volcanoes:magma}{الصهارة}؛ (ج) \omterm{def:g7:volcanoes:vent}{المنفذ}؛ (د) \omterm{def:g4:volcanoes-earthquakes-first-look:eruption}{الحمم}؛ (ه) \omterm{def:g7:volcanoes:magma-chamber}{حجرة الصهارة}.
\end{solution}

\begin{solution}{exo:g7:volcanoes:2}
(أ) انسيابي. (ب) انفجاري. (ج) انسيابي: تقذف نافورة \omterm{def:g4:volcanoes-earthquakes-first-look:eruption}{الحمم}
قطرات متوهجة تسقط ثم تجري؛ ولا تتحول \omterm{def:g7:volcanoes:magma}{الصهارة} إلى رماد.
(د) انفجاري: إنه \omterm{def:g7:volcanoes:pyroclastic-flow}{تدفق فتاتي بركاني}.
\end{solution}

\begin{solution}{exo:g7:volcanoes:3}
(أ) صواب. (ب) خطأ: أوفر الغازات البركانية هو بخار الماء؛
ويأتي ثاني أكسيد الكربون بعده. (ج) خطأ: \omterm{def:g7:volcanoes:pyroclastic-flow}{التدفق الفتاتي البركاني} انهيار
من الغاز الحار والرماد وقطع \omterm{def:g3:stones-pebbles-sand:rock}{الصخر}. (د) خطأ: إنها تقع داخل
\omterm{def:g7:plate-tectonics:plate}{صفيحة} \omterm{def:g2:water-on-earth:ocean}{المحيط} الهادئ، على بعد أكثر من 3200~كيلومتر من أقرب حد، فوق
\omterm{def:g7:volcanoes:hotspot}{بقعة ساخنة}.
\end{solution}

\begin{solution}{exo:g7:volcanoes:4}
(أ) من نحو 1000 إلى 1200~درجة. (ب) من نحو 600 إلى 800~درجة. (ج) قرب
\omterm{def:g7:volcanoes:vent}{المنفذ}، حيث \omterm{def:g4:volcanoes-earthquakes-first-look:eruption}{الحمم} صفراء (من 1000 إلى 1200~درجة) لا برتقالية
(من 800 إلى 1000~درجة): \omterm{def:g4:volcanoes-earthquakes-first-look:eruption}{فالحمم} تبرد كلما ابتعدت عن \omterm{def:g7:volcanoes:vent}{المنفذ}.
\end{solution}

\begin{solution}{exo:g7:volcanoes:5}
(أ) $838 + 221 = 1059$ \omterm{def:g4:volcanoes-earthquakes-first-look:volcano}{بركانًا}. (ب) $1059 \div 1214 \approx 0.872$:
نحو \textbf{87~\%}. (ج) $69 + 18 + 12 + 0.5 = 99.5$، أي نحو 100~\%:
والفرق الصغير ناتج عن تقريب كل نسبة مئوية.
\end{solution}

\begin{solution}{exo:g7:volcanoes:6}
(أ) $564 \div 660 \approx 0.855$: نحو \textbf{85~\%}. (ب)
$660 - 564 - 62 - 33 = 1$: مخروط واحد. (ج) المخاريط الشديدة الانحدار تبنيها
\omterm{def:g4:volcanoes-earthquakes-first-look:eruption}{الحمم} اللزجة والرماد، وهما علامة \omterm{def:g4:volcanoes-earthquakes-first-look:eruption}{الثورانات} الانفجارية: فعلى طول حزام
النار، حيث تغوص \omterm{def:g7:plate-tectonics:plate}{صفيحة} تحت أخرى، نتوقع \omterm{def:g4:volcanoes-earthquakes-first-look:eruption}{ثورانات} انفجارية كثيرة،
مثل \omterm{def:g4:volcanoes-earthquakes-first-look:eruption}{ثورانات} \omterm{def:g1:landscapes-around-us:landform}{جبل} سانت~هيلينز وبيناتوبو.
\end{solution}

\begin{solution}{exo:g7:volcanoes:7}
(أ) $1170 - 1140 = 30$~درجة. (ب) $30 \div 10 = 3$~درجات لكل
كيلومتر. (ج) في النفق لا تلامس \omterm{def:g4:volcanoes-earthquakes-first-look:eruption}{الحمم} الهواء: فالسقف
المتصلب فوقها يحبس الحرارة، كغطاء القدر. أما في
الهواء الطلق، فيفقد سطحها الحرارة في الهواء في الحال.
\end{solution}

\begin{solution}{exo:g7:volcanoes:8}
(أ) كاواي، بما لا يقل عن $5.5 - 0.7 = 4.8$~مليون سنة. (ب) كاواي: \omterm{def:g4:clouds-rain-snow:precipitation}{فالمطر}
والأمواج تبليها منذ مدة أطول بكثير. (ج) الجزيرة الكبيرة،
حيث لا تزال \omterm{def:g4:volcanoes-earthquakes-first-look:eruption}{الحمم} الجديدة تخرج. (د) نحو الشمال الغربي: فكاواي، التي حُملت
بعيدًا عن \omterm{def:g7:volcanoes:hotspot}{البقعة الساخنة} أولًا، تقع الآن شمال غربها.
\end{solution}

\begin{solution}{exo:g7:volcanoes:9}
(أ) تكاد كل \omterm{def:g4:volcanoes-earthquakes-first-look:volcano}{البراكين} تقع على \omterm{def:g7:plate-tectonics:plate-boundary}{حدود الصفائح} (87~\% في القائمة):
\omterm{def:g4:volcanoes-earthquakes-first-look:volcano}{فالبراكين} البعيدة إلى هذا الحد عن أي حد استثناء. (ب) في أعماق
\omterm{def:g6:inside-the-earth:layers}{الوشاح} تحت هاواي، تبقى \omterm{def:g7:volcanoes:hotspot}{بقعة ساخنة} حارة ملايين السنين. وهي
تصهر جزئيًا \omterm{def:g3:stones-pebbles-sand:rock}{الصخر} الذي فوقها. \omterm{def:g7:volcanoes:magma}{والصهارة}، الأخف من \omterm{def:g3:stones-pebbles-sand:rock}{الصخر} الصلب
\omterm{def:g2:water-on-earth:ocean}{المحيط} بها، تصعد عبر \omterm{def:g7:plate-tectonics:plate}{الصفيحة} وتثور. ومع تحرك \omterm{def:g7:plate-tectonics:plate}{الصفيحة}، تنمو \omterm{def:g4:volcanoes-earthquakes-first-look:volcano}{براكين}
جديدة فوق \omterm{def:g7:volcanoes:hotspot}{البقعة الساخنة} وتُحمل \omterm{def:g4:volcanoes-earthquakes-first-look:volcano}{البراكين} القديمة بعيدًا. (ج)
نقاط برتقالية («داخل \omterm{def:g7:plate-tectonics:plate}{صفيحة}»).
\end{solution}

\begin{solution}{exo:g7:volcanoes:10}
كانت \omterm{def:g7:volcanoes:magma}{الصهارة} الأولى قد فقدت معظم غازها في طريقها إلى الأعلى، كقنينة
ماء غازي تُركت مفتوحة أو واقفة: فحين بلغت السطح، لم
يبقَ ما يمزقها، فنزّت. أما \omterm{def:g7:volcanoes:magma}{صهارة} 12
و15~يونيو فكانت لا تزال تحتفظ بغازها، كالقنينة المرجوجة: فحين انخفض
الضغط قرب السطح، كوّن الغاز فقاعات دفعة واحدة،
انتفخت ومزقت \omterm{def:g7:volcanoes:magma}{الصهارة} اللزجة رمادًا. \omterm{def:g7:volcanoes:magma}{والصهارة} الأولى هي التي
كانت قد فقدت غازها.
\end{solution}

\begin{solution}{exo:g7:volcanoes:11}
(أ) $10 \div 80 = 0.125$~ساعة، و$0.125 \times 60 = 7.5$~دقيقة. (ب)
$10 \div 130 \approx 0.077$~ساعة، أي نحو $0.077 \times 60 \approx
4.6$~دقيقة، أي نحو 5~دقائق. (ج) تحتاج الأسرة إلى 20~دقيقة، لكن
التدفق يصل في 5 إلى 8~دقائق: فإذا رُئي التدفق، كان الأوان قد فات. ولا
تكون الأسرة في أمان إلا إذا رحلت حين يرى العلماء علامات
الإنذار، قبل \omterm{def:g4:volcanoes-earthquakes-first-look:eruption}{الثوران}.
\end{solution}

\begin{solution}{exo:g7:volcanoes:12}
(أ) $111 \times 360 = 39\,960$: نحو \textbf{40\,000~كيلومتر}. (ب)
$40\,000 \div 15 \approx 2667$: نحو 2\,670~كيلومترًا في اليوم. (ج)
$2667 \div 24 \approx 111$: نحو \textbf{110~كيلومترات في الساعة}. (د) الريح هي التي
تحمله: فالرماد الناعم لا سرعة خاصة له بعد أن يقذفه
\omterm{def:g4:volcanoes-earthquakes-first-look:eruption}{الثوران} إلى الأعلى؛ بل يطفو عاليًا في الهواء وتحمله الريح، وهو يسقط
ببطء.
\end{solution}

\begin{solution}{pb:g7:volcanoes:1}
\textbf{1.} من 16 إلى 31~مارس: $31 - 16 = 15$~يومًا؛ ثم 30~يومًا من
أبريل و18~يومًا من مايو: $15 + 30 + 18 = 63$~يومًا.

\textbf{2.} $10\,000 \div 62 \approx 161$: أكثر من 160 \omterm{def:g4:volcanoes-earthquakes-first-look:earthquake}{زلزالًا} صغيرًا
في اليوم، في المتوسط.

\textbf{3.} بيّن الانتفاخ أن \omterm{def:g7:volcanoes:magma}{الصهارة} صعدت عاليًا داخل
\omterm{def:g1:landscapes-around-us:landform}{الجبل} وأنها تدفع \omterm{def:g3:stones-pebbles-sand:rock}{الصخر} نحو الخارج. كما جعل السفح الشمالي
شديد الانحدار غير مستقر، مهيأً للانزلاق.

\textbf{4.} المثلث الأزرق يقع حيث تغوص \omterm{def:g7:plate-tectonics:plate}{صفيحة} تحت أخرى. ونحن
نتوقع \omterm{def:g7:volcanoes:magma}{صهارة} لزجة غنية بالغاز، \omterm{def:g4:volcanoes-earthquakes-first-look:eruption}{وثورانات} انفجارية ومخروطًا شديد الانحدار (\omterm{def:g4:volcanoes-earthquakes-first-look:volcano}{بركانًا}
طبقيًا) --- وهذا بالضبط ما تُظهره الصورة الأولى.

\textbf{5.} $240 \div 60 = 4$~كيلومترات في الدقيقة.

\textbf{6.} كان السفح المنتفخ يحبس \omterm{def:g7:volcanoes:magma}{الصهارة}، تحت الضغط،
كغطاء القنينة. فحين انزلق، انخفض الضغط في
الحال؛ فتحول الغاز الذائب في \omterm{def:g7:volcanoes:magma}{الصهارة} إلى فقاعات دفعة واحدة
واندفع خارجًا عبر السفح المفتوح --- جانبيًا، حيث زال الغطاء.

\textbf{7.} $27 \div 480 \approx 0.056$~ساعة، و$0.056 \times 60 \approx
3.4$~دقيقة. وبما أن الانفجار تحرك بهذه السرعة على الأقل، فقد استغرق على الأكثر
نحو 3.4~دقيقة.

\textbf{8.} $8 \div 130 \approx 0.062$~ساعة، و$0.062 \times 60 \approx
3.7$: نحو \textbf{4~دقائق}.

\textbf{9.} $24 \div 8.8 \approx 2.7$: ارتفع العمود نحو 2.7 مرة
أعلى من \omterm{def:g1:landscapes-around-us:landform}{جبل} إفرست.

\textbf{10.} انسيابية: خرجت \omterm{def:g4:volcanoes-earthquakes-first-look:eruption}{الحمم} اللزجة دون أن تنفجر وتكدست
فوق \omterm{def:g7:volcanoes:vent}{المنفذ}. فكالقنينة التي بقيت ساكنة، كانت هذه \omterm{def:g7:volcanoes:magma}{الصهارة} قد
فقدت معظم غازها من قبل.

\textbf{11.} قبل \omterm{def:g4:volcanoes-earthquakes-first-look:eruption}{الثوران}، كان \omterm{def:g1:landscapes-around-us:landform}{الجبل} مخروطًا أملس مدببًا مكسوًا
\omterm{def:g4:clouds-rain-snow:precipitation}{بالثلج}، فوق غابات خضراء. وبعده، زالت قمته وسفحه الشمالي: وحلّت محلهما
\omterm{def:g7:volcanoes:vent}{فوهة} هائلة مفتوحة نحونا، وصارت المنحدرات \omterm{def:g1:landscapes-around-us:landform}{والوادي}
الذي في المقدمة رمادية اللون بالرماد \omterm{def:g3:stones-pebbles-sand:rock}{والصخر} المكسَّر.

\textbf{12.} $2950 - 2549 = 401$: فقد \omterm{def:g1:landscapes-around-us:landform}{الجبل} نحو
\textbf{401~متر} من \omterm{def:g4:world-physical-map:altitude}{الارتفاع}، أي نحو 400~متر، كما رُوي في الصف~الرابع.
\end{solution}
